% ==============================================================
\section{Introduction}\label{sec:intro}
% ==============================================================

\subsection{Motivation and Application Background}

Data analysts at large retail chains want to know
``which product combinations are intensively purchased,
at which store locations, and during which time periods.''
For example, from 100 stores $\times$ 52 weeks of
point-of-sale data, the insight that
``sales of beer and potato chips are concentrated during
weeks 46--52 (year-end season) at urban stores (stores 0--9)''
directly informs region-specific inventory planning
and promotional strategies.
Similarly, in epidemiology, specific regions and periods
in which co-occurrence patterns of symptoms are concentrated
become intervention targets.
In logistics, regional and temporal surges in shipment volumes
support service optimization.

Multi-dimensional data are thus accumulated across many domains,
and extracting \emph{spatio-temporally localized dense patterns}
from such data can contribute to discovering knowledge
beyond that provided by conventional aggregate analysis.

\subsection{Limitations of Existing Methods}

In the above applications, analysts truly need to know two things:
(i)~\emph{what} is dense (itemset $P$) and
(ii)~\emph{when and where} it is dense ($d$-dimensional dense region),
i.e., the tuple $(P, R(P))$.
Pattern $P$ alone leaves ``which store to reorder to and when'' unknown,
while region $R$ alone leaves ``what to reorder'' unspecified.
Only when both are available can region- and period-specific
strategies be formulated.
However, no existing method provides both simultaneously.

\paragraph{Pattern mining methods.}
Frequent pattern mining~\cite{agrawal1994apriori,han2000mining}
relies on global support and thus misses locally dense patterns.
Existing methods that consider locality~\cite{aprioriwindow}
remain limited to a one-dimensional design on the temporal axis
and cannot exploit multi-dimensional spatial information.
\paragraph{Region extraction methods.}
Dense subtensor mining (M-Zoom~\cite{shin2016mzoom}, etc.)
cannot fit arbitrarily-shaped dense regions
due to Cartesian-product block constraints,
and it inaccurately enumerates patterns
(detailed in \S\ref{sec:related}).
Spatial statistics methods~\cite{anselin1995lisa,ord1995getis}
cannot handle the temporal dimension.
Scan statistics~\cite{kulldorff1997spatial}
lack the concept of itemsets.
ST-DBSCAN~\cite{birant2007stdbscan} relies on
point-neighborhood density and thus cannot detect
localized concentrations distinguished from background density;
it also lacks anti-monotonicity and therefore
cannot be integrated with Apriori pruning.

\subsection{Contributions of This Work}

This paper proposes
\textbf{Multi-Dimensional Apriori-Window},
which exactly extracts all dense patterns and their dense regions
over arbitrary $d$-dimensional grids
(time + $d{-}1$ spatial/categorical dimensions)
(\S\ref{sec:method}).
In a nutshell, the proposed method enables analysts to
\textbf{exactly extract ``what sells when and where''
under the sliding-window density definition}.

The technical essence of this work lies in extending
the dense interval computation of
one-dimensional Apriori-Window~\cite{aprioriwindow}
to $d$ dimensions by applying it to
the temporal multiset for each fixed spatial start point.
That is, even after multi-dimensionalization,
the core computation is essentially
a \emph{duplicate-aware one-dimensional Apriori-Window}.

Three techniques realize this multi-dimensionalization.

\begin{enumerate}
\item \textbf{Reduction to 1D per fixed spatial start point}:
For each spatial start point $\bm{s}_{\mathrm{sp}}$,
occurrences of a pattern falling into the spatial window
are collapsed into a temporal multiset
$L_P(\bm{s}_{\mathrm{sp}})$.
This decomposes the $d$-dimensional problem
into multiple one-dimensional problems
keyed by spatial start points.

\item \textbf{Apriori pruning via $(k{-}1)$-subset interval map intersection}:
Dense regions from the previous level are maintained
as temporal interval sequences per spatial start point.
For candidate patterns, these are intersected
to form candidate intervals.
This enables interval-based pruning
isomorphic to the 1D case for $k \geq 2$
(Lemma~\ref{lem:apriori}).

\item \textbf{Duplicate-aware $t$-axis stride}:
In two or more dimensions, duplicate timestamps arise.
Instead of maintaining dense-start positions one by one,
contiguous blocks of left endpoints guaranteed to be dense
by the current window alone are processed in bulk,
and output as dense intervals (Lemma~\ref{lem:stride}).
\end{enumerate}
